% =============================================================================
% APPENDIX I: Discussion Figures and Tables
% =============================================================================

\section{Discussion Figures and Tables}
\label{sec:appendix_discussion_floats}

\noindent
This section collects the figures and tables referenced in the Discussion (\cref{sec:discussion}) and Limitations (\cref{sec:limitations}) sections.

% ------------------------------------------------------------
% Toolbox Figure (from Discussion)
% ------------------------------------------------------------

\begin{figure*}[p]
\centering
\begin{tcolorbox}[
  colback=gray!3,
  colframe=black!65,
  boxrule=0.6pt,
  sharp corners,
  left=6pt,right=6pt,top=4pt,bottom=4pt,
  enhanced
]
{\small
\noindent\textbf{Toolbox: CITA objective and Conditional Safety (step-by-step).}\par
\vspace{0.35em}

\noindent\textbf{Step 1 (Preference gap).}\;
\[
\Delta_\theta(I,X;Y^+,Y^-)
=
\log \pi_\theta(Y^+\mid I,X)
-
\log \pi_\theta(Y^-\mid I,X).
\]

\noindent\textbf{Step 2 (Instruction-conditioned preference loss).}\;
\[
\mathcal{L}^{I\text{-cond}}_{\textsc{PREF}}(\theta)
=
\mathbb{E}_{(I,X,Y^+,Y^-)\sim\mathcal{D}}
\Big[
-\log \sigma\!\big(\beta\,\Delta_\theta(I,X;Y^+,Y^-)\big)
\Big].
\]

\noindent\textbf{Step 3 (Mandatory KL anchor).}\;
\[
\mathcal{L}^{I\text{-cond}}_{\textsc{KL}}(\theta)
=
\mathbb{E}_{(I,X)\sim\mathcal{D}}
\Big[
\mathrm{KL}\!\big(\pi_\theta(\cdot\mid I,X)\,\|\,\pi_0(\cdot\mid I,X)\big)
\Big].
\]

\noindent\textbf{Step 4 (Full CITA--KL).}\;
\[
\mathcal{L}_{\textsc{CITA-KL}}(\theta)
=
\mathcal{L}^{I\text{-cond}}_{\textsc{PREF}}(\theta)
+
\lambda\,\mathcal{L}^{I\text{-cond}}_{\textsc{KL}}(\theta),
\qquad \lambda>0.
\]

\vspace{0.35em}
\noindent\textbf{Step 5 (Conditional Safety as refusal-gap).}\;
\[
\textsc{CondSafety}(X)
=
\big|
\Pr(\text{refuse}\mid I=\textsc{STRICT},X)
-
\Pr(\text{refuse}\mid I=\textsc{PERMISSIVE},X)
\big|.
\]

\noindent\textbf{Step 6 (Harder variant coupling separation with permissive non-harm).}\;
\[
\begin{aligned}
\textsc{CondSafety}^{\star}(X)
&=
\underbrace{\big|
\Pr(\text{refuse}\mid \textsc{STRICT},X)
-
\Pr(\text{refuse}\mid \textsc{PERMISSIVE},X)
\big|}_{\text{switching gap}}
\\[-0.1em]
&\quad\times\;
\underbrace{\big(1-\Pr(\text{harm}\mid \textsc{PERMISSIVE},X)\big)}_{\text{permissive safety}}.
\end{aligned}
\]

\noindent\textbf{Step 7 (Distributional separation; optional diagnostic).}\;
\[
\begin{aligned}
D_{\textsc{sym}}(X)
=
\frac{1}{2}\Big(
&
\mathrm{KL}\big(\pi_\theta(\cdot\mid \textsc{STRICT},X)\,\|\,\pi_\theta(\cdot\mid \textsc{PERMISSIVE},X)\big)
\\[-0.1em]
+
&
\mathrm{KL}\big(\pi_\theta(\cdot\mid \textsc{PERMISSIVE},X)\,\|\,\pi_\theta(\cdot\mid \textsc{STRICT},X)\big)
\Big).
\end{aligned}
\]
}
\end{tcolorbox}
\caption{\textbf{Toolbox: large equations moved out of the two-column flow.} We present the CITA--KL objective and Conditional Safety diagnostics as short, column-safe steps. In the main text we reference this toolbox and keep the discussion narrative tight.}
\label{fig:toolbox_cita_condsafety}
\end{figure*}

% ------------------------------------------------------------
% Limitations Table (from Limitations section)
% ------------------------------------------------------------

\begin{table*}[ht!]
\centering
\normalsize
\setlength{\tabcolsep}{4pt}
\renewcommand{\arraystretch}{1.2}
\begin{tabularx}{\textwidth}{@{}l>{\raggedright\arraybackslash}X@{}}
\toprule
\textbf{Limitation} & \textbf{Impact \& Mitigation Path} \\
\midrule
Single model (Llama-3.1-8B) & Validate on GPT, Mistral, Gemma; scale to 70B \\ \hline
English-only evaluation & Extend to multilingual instruction transfer \\ \hline
10 instruction types & Add domain-specific types (medical, legal) \\ \hline
Single dataset (PKU-SafeRLHF) & Test on diverse preference datasets \\
\bottomrule
\end{tabularx}
\caption{Current experimental limitations and mitigation paths.}
\label{tab:limitations}
\end{table*}

% ------------------------------------------------------------
% Discussion Claims Table (from D.1)
% ------------------------------------------------------------

\begin{table*}[ht!]
\centering
\normalsize
\setlength{\tabcolsep}{6pt}
\adjustbox{max width=\textwidth}{%
\begin{tabular}{p{3.2cm}p{6.9cm}p{5.9cm}}
\toprule
\textbf{Discussion claim} & \textbf{Operational meaning (what must be true)} & \textbf{Primary substantiation in this paper} \\
\midrule
\textbf{C1: Switching as a control primitive} &
\textbf{Same $X$} should yield \textbf{distinct behavioral contracts} under different $I$; switching must be \textbf{stable} (no leakage/oscillation/unsafe drift). &
Prompt-held-constant isolation + reliability framing in ECLIPTICA; prompt-matched paired comparisons. \\
\textbf{C2: Mandatory KL is structural} &
Training must preserve a \textbf{family} $\{\pi_\theta(\cdot\mid I,\cdot)\}$ of \textbf{nearby instruction-indexed optima} rather than collapsing to an instruction-agnostic shortcut. &
Mandatory KL as a trust-region anchor; geometric (locality/Fisher) rationale; objective summarized in Fig.~\ref{fig:toolbox_cita_condsafety} (Steps~1--4). \\
\textbf{C3: Strong instruction sensitivity} &
Adding $I$ should change behavior (relative to NoInstruct) on \textbf{controlled-switching probes}, not merely improve surface ``instruction awareness.'' &
$\Delta=\textsc{Instruct}-\textsc{NoInstruct}$ protocol and the \cref{tab:improvement_comparison} sensitivity pattern across probes. \\
\bottomrule
\end{tabular}%
}
\vspace{-0.3em}
\caption{\textbf{Discussion map.} Claims, their operational content, and where the paper substantiates them.}
\label{tab:discussion_claims}
\vspace{-0.9em}
\end{table*}

% ------------------------------------------------------------
% Benchmark Tradeoffs Table (from D.3)
% ------------------------------------------------------------

\begin{table*}[ht!]
\centering
\normalsize
\setlength{\tabcolsep}{6pt}
\adjustbox{max width=\textwidth}{%
\begin{tabular}{p{3.0cm}p{5.2cm}p{7.1cm}}
\toprule
\textbf{Probe} & \textbf{What it rewards} & \textbf{Failure modes it surfaces (switching lens)} \\
\midrule
ECLIPTICA & \textbf{Instruction awareness} (fidelity $\times$ shift). & High via large stylistic shifts; may not ensure \textbf{neighborhood-stable regimes}. \\
TruthfulQA (HON--CONF) & \textbf{Calibration} separation under uncertainty. & Overconfident collapse; instruction-insensitive calibration. \\
Cond.\ Safety (STRICT/PERMISSIVE) & \textbf{Binary} refusal-gap. & Sharp boundary without guaranteeing permissive-mode quality. \\
Length Control (DETAIL/CONCISE) & \textbf{Explicit} verbosity contract. & Regime leakage; verbosity collapse under competing constraints. \\
AQI & \textbf{Axiom-level} clustering structure. & Collapse into a single policy basin; incoherent axiom mixing. \\
\bottomrule
\end{tabular}%
}
\vspace{-0.3em}
\caption{\textbf{Benchmark patterns as diagnostics.} Each probe emphasizes a distinct aspect of instruction-conditioned control; the observed method ordering is best read as a \textbf{multi-axis switching profile}, not a single scalar ``best.''}
\label{tab:tradeoffs}
\vspace{-0.9em}
\end{table*}

% ------------------------------------------------------------
% Discussion & Limitations At-a-Glance Table
% ------------------------------------------------------------

% --- icon helpers ---
\newcommand{\icnUse}{\faCompass}              % usage / how-to
\newcommand{\icnScope}{\faBullseye}           % scope / what it captures
\newcommand{\icnWarn}{\faExclamationTriangle} % caution / sensitivity
\newcommand{\icnFix}{\faFlask}                % experiment / fix
\newcommand{\icnDont}{\faBan}                 % do-not-use-as-gate
\newcommand{\icnRel}{\faClipboardList}        % reproducibility checklist
\newcommand{\icnRoad}{\faRoute}               % roadmap
\newcommand{\icnTool}{\faToolbox}             % toolbox figure pointer
\newcommand{\icnSwap}{\faExchange*}           % works in more setups
\newcommand{\icnGeom}{\faProjectDiagram}      % geometry / manifold
\newcommand{\icnShield}{\faShieldAlt}         % safety boundary
\newcommand{\icnChart}{\faChartLine}          % pattern / metrics

% --- column type (C only; L already defined) ---
\newcolumntype{C}[1]{>{\centering\arraybackslash}p{#1}}

% --- Table 1: Discussion ---
\begin{table*}[ht!]
\centering
\adjustbox{max width=\textwidth}{%
\normalsize
\renewcommand{\arraystretch}{1.18}
\setlength{\tabcolsep}{4.2pt}
\begin{tabular}{L{0.115\textwidth} L{0.23\textwidth} L{0.295\textwidth} L{0.30\textwidth}}
\toprule
\textbf{Block} &
\textbf{\icnScope\ \,What it is for (read-out)} &
\textbf{\icnWarn\ \,What to watch (failure / sensitivity)} &
\textbf{\icnFix\ \,What fixes it (report / experiment)} \\
\midrule

\rowcolor{gray!10}
\multicolumn{4}{l}{\textbf{Discussion (how to read ECLIPTICA + \textsc{CITA})}}\\[-0.5mm]
\rowcolor{gray!10}\multicolumn{4}{l}{\rule{0pt}{2.6ex}}\\[-2.2mm]

\textbf{D1 \icnUse\ \,Causal isolation} &
\textbf{Prompt-held-constant} switching: attribute changes to $I$ (not $X$); paired deltas are the unit of evidence. &
Leakage via prompt artifacts or instruction templates: apparent "switching" could be superficial style drift. &
Always report \textbf{paired} comparisons ($\Delta=\textsc{Instruct}-\textsc{NoInstruct}$); include template ablations / paraphrased $I$ sanity checks. \\

\textbf{D2 \icnSwap\ \,Switching primitive} &
\textbf{One checkpoint, many regimes}: distinct behavioral contracts for the same $X$ under different $I$; stable bidirectional toggling. &
Mode interference: regimes collapse into a single implicit behavior; oscillation across instructions; "leaky" mixtures. &
Use objective + diagnostics that preserve a \textbf{policy family} $\{\pi_\theta(\cdot\mid I,\cdot)\}$; emphasize stability checks under repeated toggles. \\

\textbf{D3 \icnGeom\ \,Why mandatory KL} &
KL provides \textbf{geometry}: encourages \textbf{nearby optima} so switching is local, controlled, and non-collapsing. &
Without anchoring, preference gradients across many $I$ can produce instruction-agnostic shortcuts or brittle over-steering. &
Reference \textbf{Toolbox Fig.~\ref{fig:toolbox_cita_condsafety}} (Steps~1--4); report $\lambda$ sensitivity and regime-collapse indicators. \\

\textbf{D4 \icnChart\ \,Profile, not scalar} &
Leaderboard is a \textbf{multi-axis switching profile}: some probes test "does it move," others test "does structure persist." &
Over-reading a single metric (e.g., ECLIPTICA or Cond.\ Safety) as "best alignment" misstates what is being certified. &
Always pair \textbf{discrete} boundary probes with \textbf{continuous} regime probes (calibration, length, axiom structure); analyze cross-metric disagreements. \\

\bottomrule
\end{tabular}%
}
\caption{\textbf{Discussion at a glance.} How to read ECLIPTICA and \textsc{CITA}: what the benchmark certifies (prompt-held-constant switching) and where the method's geometry enters (mandatory KL; see Toolbox Fig.~\ref{fig:toolbox_cita_condsafety}).}
\label{tab:ecliptica_discussion_glance}
\end{table*}

% --- Table 2: Limitations ---
\begin{table*}[ht!]
\centering
\adjustbox{max width=\textwidth}{%
\normalsize
\renewcommand{\arraystretch}{1.18}
\setlength{\tabcolsep}{4.2pt}
\begin{tabular}{L{0.115\textwidth} L{0.23\textwidth} L{0.295\textwidth} L{0.30\textwidth}}
\toprule
\textbf{Block} &
\textbf{\icnScope\ \,What it is for (read-out)} &
\textbf{\icnWarn\ \,What to watch (failure / sensitivity)} &
\textbf{\icnFix\ \,What fixes it (report / experiment)} \\
\midrule

\rowcolor{gray!10}
\multicolumn{4}{l}{\textbf{Limitations (what can break and why it matters)}}\\[-0.5mm]
\rowcolor{gray!10}\multicolumn{4}{l}{\rule{0pt}{2.6ex}}\\[-2.2mm]

\textbf{L1 \icnWarn\ \,Instruction semantics} &
Switching depends on $I$ being semantically separable and consistently interpreted by the model. &
Paraphrase sensitivity: small changes in $I$ wording may change regime identity or induce partial mixing. &
Report \textbf{instruction paraphrase} robustness; include a minimal "instruction invariance" panel (synonyms/format perturbations). \\

\textbf{L2 \icnWarn\ \,Objective specificity} &
Claims are tied to instruction-conditioned preference learning with \textbf{mandatory} KL anchoring. &
Generalization to RLHF/constitutional or tool-augmented stacks may shift where and how regimes form. &
Run matched comparisons across objectives (DPO/GRPO/RLHF variants) under the same ECLIPTICA grid; keep within-family deltas. \\

\textbf{L3 \icnWarn\ \,Judge drift / measurement} &
Metrics rely on fixed checkers and scoring conventions. &
Binary refusal or harm checkers can drift; calibration probes can be sensitive to prompt framing and evaluator policy. &
Version all judges; report stability across judge variants; add a small robustness appendix (prompt subsets, seed/decoder checks). \\

\textbf{L4 \icnWarn\ \,Cond.\ Safety degeneracy} &
Cond.\ Safety is a \textbf{mode-separation} diagnostic (STRICT vs PERMISSIVE). &
High refusal-gap can be achieved by degenerate extremes ("refuse always" vs "comply always") if permissive quality is unchecked. &
Use the \textbf{harder} definition in Toolbox Fig.~\ref{fig:toolbox_cita_condsafety} (Step~6); add permissive-mode harm/quality checks. \\

\textbf{L5 \icnWarn\ \,Distributional mismatch} &
ECLIPTICA certifies controlled switching on its prompt-instruction grid. &
Deployment prompts may differ (domain shift, adversarial phrasing); regimes may not transfer cleanly. &
Stratify by prompt category; add out-of-grid evaluation (held-out domains, adversarial paraphrases, long-context variants). \\

\textbf{L6 \icnDont\ \,Not a deployment gate} &
Useful as \textbf{control-channel} evidence and switching diagnostics; complements behavioral suites. &
A single score cannot certify safety or truthfulness under all conditions; disagreement cases are informative, not "noise." &
State explicitly: \textbf{not a pass/fail gate}; use disagreements to drive targeted eval and causal analysis. \\

\bottomrule
\end{tabular}%
}
\caption{\textbf{Limitations at a glance.} What can break and why it matters: each limitation, its failure mode, and the recommended check or experiment.}
\label{tab:ecliptica_limitations_glance}
\end{table*}

% --- Table 3: Roadmap ---
\begin{table*}[ht!]
\centering
\adjustbox{max width=\textwidth}{%
\normalsize
\renewcommand{\arraystretch}{1.18}
\setlength{\tabcolsep}{4.2pt}
\begin{tabular}{L{0.115\textwidth} L{0.23\textwidth} L{0.295\textwidth} L{0.30\textwidth}}
\toprule
\textbf{Block} &
\textbf{\icnScope\ \,What it is for (read-out)} &
\textbf{\icnWarn\ \,What to watch (failure / sensitivity)} &
\textbf{\icnFix\ \,What fixes it (report / experiment)} \\
\midrule

\rowcolor{gray!10}
\multicolumn{4}{l}{\textbf{Roadmap (high-level, testable directions)}}\\[-0.5mm]
\rowcolor{gray!10}\multicolumn{4}{l}{\rule{0pt}{2.6ex}}\\[-2.2mm]

\textbf{FW \icnRoad\ \,Next steps} &
Turn switching into an auditable standard: benchmark + objective + diagnostics for instruction-conditioned regimes. &
Overcommitting speculation; roadmap should remain crisp, measurable, and reproducible. &
(1) Extend ECLIPTICA across architectures/objectives; (2) add causal tests for regime locality; (3) publish standardized prompt+instruction packs + judge versioning + robustness panel. \\

\bottomrule
\end{tabular}%
}
\caption{\textbf{Roadmap at a glance.} High-level, testable future directions for instruction-conditioned alignment.}
\label{tab:ecliptica_roadmap_glance}
\end{table*}
